% .claude/skills/md-to-pdf/assets/templates/resume/main.tex — レジュメの雛形
% 発表・ゼミ・授業の配布資料用（LuaLaTeX + jlreq）
%
% 使い方:
%   1. このフォルダごと work/<プロジェクト>/drafts/<名前>/ にコピーする
%   2. 書き換えるのはタイトル情報と本文だけ
%   3. latexmk main.tex  （VSCode の LaTeX Workshop なら保存時に自動ビルド）
%
% 参考文献は biblatex ではなく thebibliography（手書き）を使う。
% 数点しか挙げないレジュメでは .bib を用意するより速い。

\documentclass[a4paper,11pt,article]{jlreq}

% ===== フォント =====
\usepackage{luatexja-fontspec}
% 和文明朝: 原ノ味明朝（TeX Live 2020以降に同梱）→ Noto → IPAex明朝
\IfFontExistsTF{HaranoAjiMincho}{%
  \setmainjfont{HaranoAjiMincho}%
}{%
  \IfFontExistsTF{Noto Serif CJK JP}{\setmainjfont{Noto Serif CJK JP}}{\setmainjfont{IPAexMincho}}%
}
% 和文ゴシック: macOS はヒラギノ、Linux/Windows は Noto → 原ノ味 → IPAex
\IfFontExistsTF{Hiragino Kaku Gothic ProN}{%
  \setsansjfont{Hiragino Kaku Gothic ProN}[
    UprightFont = {Hiragino Kaku Gothic ProN W3},
    BoldFont    = {Hiragino Kaku Gothic ProN W6},
  ]%
}{%
  \IfFontExistsTF{Noto Sans CJK JP}{%
    \setsansjfont{Noto Sans CJK JP}%
  }{%
    \IfFontExistsTF{HaranoAjiGothic}{\setsansjfont{HaranoAjiGothic}}{\setsansjfont{IPAexGothic}}%
  }%
}
% 欧文: TeX Gyre Termes（Times系、TeX Live同梱）
\setmainfont{TeX Gyre Termes}

%%% 古い TeX Live への保険 %%%
% \AddToHook は LaTeX 2020-10 以降のカーネル機能。それ以前の TeX Live では
% 未定義でビルドが止まる。見出しのゴシック化を諦めてでも通るようにしておく
% （現行の TeX Live では下の行は何もしない）。
\providecommand{\AddToHook}[2]{}

% ===== 見出しをゴシック体に =====
\AddToHook{begindocument/before}{%
  \ModifyHeading{section}{font={\sffamily\bfseries\Large}}%
  \ModifyHeading{subsection}{font={\sffamily\bfseries\large}}%
  \ModifyHeading{subsubsection}{font={\sffamily\bfseries\normalsize}}%
}
% 章番号を出したくないときは secnumdepth=0
\setcounter{secnumdepth}{3}
\setcounter{tocdepth}{3}

% ===== 紙面・行間 =====
\usepackage[
  top=25mm, bottom=25mm,
  left=25mm, right=20mm,
]{geometry}
% 行間は jlreq デフォルトを使用（\onehalfspacing は使わない）
% さらに詰めたい場合: \linespread{0.95}\selectfont を本文側で指定

% ===== 箇条書き（多段インデント） =====
% leftmargin の値（1.5em）を変えるとインデント幅を一括で変更できる
% ラベルを別記号に変えたい場合: label={\textbullet} や label={\textendash} に
\usepackage{enumitem}
\setlist[itemize]{nosep, topsep=0.3em, partopsep=0pt}
\setlist[itemize,1]{label={・}, leftmargin=1.5em}
\setlist[itemize,2]{label={・}, leftmargin=1.5em}
\setlist[itemize,3]{label={・}, leftmargin=1.5em}
\setlist[itemize,4]{label={・}, leftmargin=1.5em}

% ===== 図表 =====
\usepackage{graphicx}
\usepackage{wrapfig}    % 小さい図を本文に回り込ませる
\usepackage{booktabs}
\usepackage{caption}
\captionsetup{font=small, labelfont=bf}

% ===== 色・ハイパーリンク =====
\usepackage[dvipsnames,svgnames,x11names]{xcolor}
\usepackage{hyperref}
\hypersetup{
  colorlinks=true, linkcolor=black, citecolor=black,
  urlcolor=blue!60!black, bookmarksnumbered=true,
}

% ===== 出典マクロ：[著者 年] / [著者 年, pp.X-Y] / [著者 年; 著者 年] =====
% 中身をそのまま角括弧で囲うだけのシンプル実装。
% 例: \citep{山田 2020}
%     \citep{鈴木 2023, pp.10--13}
%     \citep{山田 2020; 鈴木 2023, pp.10--13}
\newcommand{\citep}[1]{[#1]}

% ===== テーゼ用囲み枠 =====
\usepackage{tcolorbox}
\tcbuselibrary{breakable, skins}
\newtcolorbox{thesis}{
  enhanced, breakable, colback=white, colframe=black,
  boxrule=0.8pt, arc=0pt,
  left=8pt, right=8pt, top=6pt, bottom=6pt,
  before skip=0.8em, after skip=0.8em,
}

% ===== 参考文献リストの余白調整 =====
% 標準の thebibliography は右余白が広めに残るので、左右の余白を詰めて
% 本文と同じ幅まで使うように再定義する。
\makeatletter
\renewenvironment{thebibliography}[1]{%
  \section*{\refname}%
  \@mkboth{\MakeUppercase\refname}{\MakeUppercase\refname}%
  \list{[\@arabic\c@enumiv]}{%
    \settowidth\labelwidth{[#1]}%
    \leftmargin\labelwidth
    \advance\leftmargin\labelsep
    \setlength{\rightmargin}{0pt}%
    \setlength{\itemsep}{0.3em}%
    \setlength{\parsep}{0pt}%
    \@openbib@code
    \usecounter{enumiv}%
    \let\p@enumiv\@empty
    \renewcommand\theenumiv{\@arabic\c@enumiv}%
  }%
  \sloppy
  \clubpenalty4000
  \@clubpenalty\clubpenalty
  \widowpenalty4000%
  \sfcode`\.\@m
}{%
  \def\@noitemerr{\@latex@warning{Empty `thebibliography' environment}}%
  \endlist
}
\makeatother

% ===== ページ番号 =====
\usepackage{lastpage}
\usepackage{fancyhdr}
\pagestyle{fancy}
\fancyhf{}
\renewcommand{\headrulewidth}{0pt}
\fancyfoot[R]{\thepage\,/\,\pageref{LastPage}}

% ===== タイトル書式（サブタイトル＋所属右寄せ対応） =====
\makeatletter
\newif\if@hassubtitle    \@hassubtitlefalse
\newif\if@hasaffiliation \@hasaffiliationfalse
\newcommand{\subtitle}[1]{\@hassubtitletrue\def\@subtitle{#1}}
\newcommand{\affiliation}[1]{\@hasaffiliationtrue\def\@affiliation{#1}}

\renewcommand{\maketitle}{%
  \begin{center}
    {\LARGE\sffamily\bfseries \@title}\par
    \if@hassubtitle
      \vspace{0.4em}{\large\sffamily \@subtitle}\par
    \fi
  \end{center}
  \begin{flushright}
    \if@hasaffiliation \@affiliation\par \fi
    \ifx\@author\@empty\else \@author\par \fi
    \ifx\@date\@empty\else   \@date\par   \fi
  \end{flushright}
  \vspace{0.5em}%
}
\makeatother

% ===== タイトル情報 =====
\title{レジュメのタイトル}
\subtitle{サブタイトルをここに置く}
\author{氏名}
\affiliation{所属（大学・研究科名）}
\date{\today}

\begin{document}
\maketitle

% =========================================================
\section{はじめに}

\begin{itemize}
  \item 出典の書き方
    \begin{itemize}
      \item 単一文献: \citep{山田 2020}
      \item ページ指定: \citep{鈴木 2023, pp.10--13}
      \item 複数文献: \citep{山田 2020; 鈴木 2023, pp.10--13}
    \end{itemize}
  \item 多段インデント（\verb|leftmargin=1.5em| を調整してインデント幅を変更可）
    \begin{itemize}
      \item 第2階層
        \begin{itemize}
          \item 第3階層
            \begin{itemize}
              \item 第4階層
            \end{itemize}
        \end{itemize}
    \end{itemize}
\end{itemize}

\begin{thesis}
中心となる主張（テーゼ）はこの囲み枠で強調する。
\end{thesis}

% =========================================================
\section{本論}

\subsection{大きな図：中央寄せ}

\begin{figure}[h]
  \centering
  \includegraphics[width=0.7\linewidth]{example-image-a}
  \caption{大きい図は \texttt{figure} 環境＋\texttt{centering} で中央寄せ}
  \label{fig:big}
\end{figure}

\begin{itemize}
  \item 紙面幅に余裕がある図はこちらで全幅を確保して配置
\end{itemize}

\subsection{小さな図：本文の右に回り込ませる}

% wrapfigure の使い方:
%  - 第1オプション [N] は図が占有する行数を指定する（並走テキストが短くて
%    も次ページに溢れない）。図の高さに合わせて 8〜14 程度を調整する。
%  - 並走テキストが足りないと次ページの参考文献などに食い込むので、
%    必ず行数指定 [N] を入れるのが安全。
\begin{wrapfigure}[10]{r}{0.35\linewidth}
  \centering
  \includegraphics[width=0.95\linewidth]{example-image-b}
  \caption{\texttt{wrapfigure} で本文右に配置}
  \label{fig:small}
\end{wrapfigure}

\begin{itemize}
  \item 図の幅を \verb|{r}{0.35\linewidth}| で指定（\texttt{r}=右、\texttt{l}=左）
  \item 行数を \verb|[10]| で固定すると次ページに溢れない
  \item 図表の脇に置きたい注釈や箇条書きに便利
  \item 引用例: \citep{Smith 2019, pp.20--25}
  \item さらに本文を加えて、図と本文の高さを揃える
  \item 必要に応じてもう一段下げる
    \begin{itemize}
      \item サブ項目を並べる
      \item 出典付き \citep{Doe 2021}
    \end{itemize}
\end{itemize}

\subsection{表}

\begin{table}[h]
  \centering
  \caption{表の例}
  \begin{tabular}{lll}
    \toprule
    項目 & 前期 & 後期 \\
    \midrule
    件数 & 70件  & 100件 \\
    区分数 & 6区分 & 8区分 \\
    \bottomrule
  \end{tabular}
\end{table}

% =========================================================
\section{おわりに}

\begin{itemize}
  \item 要点を再掲
  \item 残された課題 \citep{鈴木 2023, pp.30--31}
\end{itemize}

% =========================================================
% \clearpage で wrapfigure などのフロート影響を完全に終わらせてから
% 参考文献に入る（紙面が短い場合は \par\bigskip に置換しても可）
\clearpage
\renewcommand{\refname}{参考文献}
\begin{thebibliography}{99}
  \bibitem{yamada2020}
    山田太郎（2020）『サンプル書籍のタイトル――副題の表示確認』学術出版社.
  \bibitem{suzuki2023}
    鈴木花子（2023）「論文題目」『雑誌名』第××号、pp.10--13.
  \bibitem{smith2019}
    Smith, John (2019) \emph{A Sample Monograph Title}, Cambridge: Cambridge University Press.
  \bibitem{doe2021}
    Doe, Jane (2021) ``A Sample Article Title,'' \emph{Journal of Sample Studies}, 12(3), pp.35--58.
\end{thebibliography}

\end{document}
